\documentclass[../../ddppaper2.tex]{subfiles}

\makeatletter
\@ifundefined{preamble@loaded}{
    \def\preamble@loaded{}

\input{../common/packages}
\input{../common/definitions}
\input{../common/abbreviations}
\makenoidxglossaries

\makeatletter  % set paper size if not included in presentation
\@ifclassloaded{beamer}{
    \usepackage[
        backend=biber,
        style=numeric,
        sorting=none
    ]{biblatex}
    \addbibresource{../common/bibliography/bibliography.bib}
}{
    \ifdefined\paper
        % Apply this only when compiling the main paper
        \geometry{letterpaper}
    \fi

    \ifdefined\thesis
        % Apply this only when compiling the thesis
        \geometry{a4paper}
    \fi

}
\makeatother

% Package additions from REVTeX templates
\usepackage{dcolumn}% Align table columns on decimal point
\usepackage{hyperref}% Load after packages that redefine cross-reference commands

\ifdefined\paper
    \usepackage[mathlines]{lineno}
    \linenumbers\relax % Commence numbering lines in the paper
\fi
\ifdefined\thesis
    \usepackage[mathlines]{lineno}
    \linenumbers\relax % Commence numbering lines in the thesis
\fi

% Initialization of tikz lib
\usetikzlibrary{fit,calc,shadows}

% Optional geometry test block from apssamp template
%\usepackage[showframe,%Uncomment any one of the following lines to test
%%scale=0.7, marginratio={1:1, 2:3}, ignoreall,% default settings
%%text={7in,10in},centering,
%%margin=1.5in,
%%total={6.5in,8.75in}, top=1.2in, left=0.9in, includefoot,
%%height=10in,a5paper,hmargin={3cm,0.8in},
%]{geometry}

}{}
\makeatother


\begin{document}

\subsubsection{Background Only Distribution}

Wilks and Wald's theorems gives us an expected behavior of the test statistic we constructed. According to Wilks~\cite{Wilks_Original} and Wald~\cite{Wald_Original}, we have that the distribution of $\invllratio$ itself should converge to a $\chi^2_d$

\begin{equation}
    \invllratio \sim \chi^2_{d}
\end{equation}

\noindent with $d$ of the number of degrees of freedom on our signal hypothesis minus those of the null hypothesis.

and so, given no new signal in \gls{sr}, we should count the number of free parameters in our signal hypothesis and be able to approximate the behavior of $\invllratio$.

\subsubsection{In the Lack of Pure Background}

When examining a real experimental sample, we would not be able to extract properties from a background-only toy, since all we have is the sample itself that may or may not contain a signal.
As in the previous work~\cite{LastDDPPaper}, we observe this background-only behavior by performing permutation tests - shuffling the $\dataset{A}$ and $\dataset{B}$ samples, then estimating the test statistic. Repeating this process generates a distribution that is indistinguishable from background only behavior, as demonstrated in Figure~\ref{fig:permutation_test}

\begin{figure}
    \centering
    \includegraphics[width=\textwidth, trim=0 12bp 0 0, clip]{}
    \caption{[PLACEHOLDER] Permutation tests performed on a sample from CMS Open data \cite{CMS_OpenData_VVTo2L2Nu_2016}.}
    \label{fig:permutation_test}
\end{figure}

\end{document}